% !TEX TS-program = XeLaTeX
% use the following command:
% all document files must be coded in UTF-8
\documentclass[portuguese]{textolivre}
% build HTML with: make4ht -e build.lua -c textolivre.cfg -x -u article "fn-in,svg,pic-align"

\journalname{Texto Livre}
\thevolume{19}
%\thenumber{1} % old template
\theyear{2026}
\receiveddate{\DTMdisplaydate{2025}{12}{20}{-1}} % YYYY MM DD
\accepteddate{\DTMdisplaydate{2026}{4}{1}{-1}}
\publisheddate{\DTMdisplaydate{2026}{8}{14}{-1}}
\corrauthor{João Antonio Lemes Bueno}
\articledoi{10.1590/1983-3652.2026.63582}
%\articleid{NNNN} % if the article ID is not the last 5 numbers of its DOI, provide it using \articleid{} commmand 
% list of available sesscions in the journal: articles, dossier, reports, essays, reviews, interviews, editorial
\articlesessionname{articles}
\runningauthor{Bueno e Terçariol} 
%\editorname{Leonardo Araújo} % old template
\sectioneditorname{Daniervelin Pereira~\orcid{0000-0003-1861-3609}}
\layouteditorname{Rayany Chaves~\orcid{0009-0007-0028-484X}}

\title{Evidências qualitativas do uso de videogame no Ensino Médio Integrado ao Técnico: uma abordagem comportamental para a Educação Baseada em Evidências}
\othertitle{Qualitative evidence of video game use in Technical Integrated High School: a behavioral approach to Evidence-Based Education}
% if there is a third language title, add here:
%\othertitle{Artikelvorlage zur Einreichung beim Texto Livre Journal}

\author[1]{João Antonio Lemes Bueno~\orcid{0009-0002-4073-8934}\thanks{Email: \href{mailto:jl.bueno@uni9.edu.br}{jl.bueno@uni9.edu.br}}}
\author[1]{Adriana Aparecida de Lima Terçariol~\orcid{0000-0002-5824-2294}\thanks{Email: \href{mailto:atercariol@gmail.com}{atercariol@gmail.com}}}
\affil[1]{Universidade Nove de Julho, Programa de Pós-Graduação em Educação, SP, Brasil.}



\addbibresource{article.bib}
% use biber instead of bibtex
% $ biber article

% used to create dummy text for the template file
\definecolor{dark-gray}{gray}{0.35} % color used to display dummy texts
\usepackage{lipsum}
\SetLipsumParListSurrounders{\colorlet{oldcolor}{.}\color{dark-gray}}{\color{oldcolor}}

% used here only to provide the XeLaTeX and BibTeX logos
\usepackage{hologo}

% if you use multirows in a table, include the multirow package
\usepackage{multirow}

% provides sidewaysfigure environment
\usepackage{rotating}

% CUSTOM EPIGRAPH - BEGIN 
%%% https://tex.stackexchange.com/questions/193178/specific-epigraph-style
\usepackage{epigraph}
\renewcommand\textflush{flushright}
\makeatletter
\newlength\epitextskip
\pretocmd{\@epitext}{\em}{}{}
\apptocmd{\@epitext}{\em}{}{}
\patchcmd{\epigraph}{\@epitext{#1}\\}{\@epitext{#1}\\[\epitextskip]}{}{}
\makeatother
\setlength\epigraphrule{0pt}
\setlength\epitextskip{0.5ex}
\setlength\epigraphwidth{.7\textwidth}
% CUSTOM EPIGRAPH - END

% to use IPA symbols in unicode add
%\usepackage{fontspec}
%\newfontfamily\ipafont{CMU Serif}
%\newcommand{\ipa}[1]{{\ipafont #1}}
% and in the text you may use the \ipa{...} command passing the symbols in unicode

% LANGUAGE - BEGIN
% ARABIC
% for languages that use special fonts, you must provide the typeface that will be used
% \setotherlanguage{arabic}
% \newfontfamily\arabicfont[Script=Arabic]{Amiri}
% \newfontfamily\arabicfontsf[Script=Arabic]{Amiri}
% \newfontfamily\arabicfonttt[Script=Arabic]{Amiri}
%
% in the article, to add arabic text use: \textlang{arabic}{ ... }
%
% RUSSIAN
% for russian text we also need to define fonts with support for Cyrillic script
% \usepackage{fontspec}
% \setotherlanguage{russian}
% \newfontfamily\cyrillicfont{Times New Roman}
% \newfontfamily\cyrillicfontsf{Times New Roman}[Script=Cyrillic]
% \newfontfamily\cyrillicfonttt{Times New Roman}[Script=Cyrillic]
%
% in the text use \begin{russian} ... \end{russian}
% LANGUAGE - END

% EMOJIS - BEGIN
% to use emoticons in your manuscript
% https://stackoverflow.com/questions/190145/how-to-insert-emoticons-in-latex/57076064
% using font Symbola, which has full support
% the font may be downloaded at:
% https://dn-works.com/ufas/
% add to preamble:
% \newfontfamily\Symbola{Symbola}
% in the text use:
% {\Symbola }
% EMOJIS - END

% LABEL REFERENCE TO DESCRIPTIVE LIST - BEGIN
% reference itens in a descriptive list using their labels instead of numbers
% insert the code below in the preambule:
%\makeatletter
%\let\orgdescriptionlabel\descriptionlabel
%\renewcommand*{\descriptionlabel}[1]{%
%  \let\orglabel\label
%  \let\label\@gobble
%  \phantomsection
%  \edef\@currentlabel{#1\unskip}%
%  \let\label\orglabel
%  \orgdescriptionlabel{#1}%
%}
%\makeatother
%
% in your document, use as illustraded here:
%\begin{description}
%  \item[first\label{itm1}] this is only an example;
%  % ...  add more items
%\end{description}
% LABEL REFERENCE TO DESCRIPTIVE LIST - END


% add line numbers for submission
%\usepackage{lineno}
%\linenumbers

\begin{document}
\maketitle

\begin{polyabstract}
\begin{abstract}
Este artigo analisa, sob uma ótica analítico-comportamental e da Educação Baseada em Evidências (EBE), os efeitos do uso do videogame \textit{Lenin – The Lion} \cite{bueno2019lenin} em estudantes do Ensino Médio Integrado ao Técnico (EMIT). A investigação qualitativa utilizou observação participante, diário de bordo e rodas de conversa para a codificação funcional de comportamentos. Os resultados indicam que a narrativa do jogo atua como antecedente discriminativo para a apropriação de conceitos e construção de conhecimento declarativo, intensificando trocas colaborativas (mandos e intraverbais) e estimulando repertórios de auto-observação e o engajamento sustentado. Conclui-se que o videogame, mediante mediação pedagógica, opera como ferramenta eficaz para a aprendizagem e o desenvolvimento de competências socioemocionais ao criar um ambiente seguro para simulação ética. O estudo sugere a manualização de práticas e protocolos de observação para fortalecer a implementação de jogos na EBE.

\keywords{Videogames \sep Ensino Médio \sep Prática pedagógica \sep  Análise do comportamento \sep Comportamento verbal} 
\end{abstract}

\begin{english} 
\begin{abstract}
This article analyzes, from a behavior-analytic perspective and within the framework of Evidence-Based Education (EBE), the effects of using the video game \textit{Lenin – The Lion} \cite{bueno2019lenin} on students of an Integrated Technical High School (EMIT). The qualitative investigation employed participant observation, field notes, and conversation circles for the functional coding of behaviors. Results indicate that the game narrative acts as a discriminative antecedent for concept appropriation and the construction of declarative knowledge, intensifying collaborative exchanges (mands and intraverbals) and stimulating self-observation repertoires and sustained engagement. It is concluded that the video game, through pedagogical mediation, operates as an effective tool for learning and the development of social-emotional competencies by creating a safe environment for ethical simulation. The study suggests the manualization of practices and observation protocols to strengthen the implementation of games in EBE.

\keywords{Video games \sep Secondary education \sep Teaching practice \sep Behavior analysis \sep Verbal behavior}
\end{abstract}
\end{english}
% if there is another abstract, insert it here using the same scheme
\end{polyabstract}

\section{Introdução}\label{sec-intro}
A Educação Baseada em Evidências (EBE) configura-se contemporaneamente como um processo de tomada de decisão pedagógica, orientado por um imperativo ético e profissional. Tal processo exige que as escolhas educacionais não sejam arbitrárias, sustentando-as por um tripé fundamental: as melhores evidências científicas disponíveis, a expertise do profissional da educação e a consideração rigorosa das características e necessidades do contexto local e dos estudantes.

A literatura especializada \cite{davies1999evidence,slavin1996education,young2016conhecimento} enfatiza que a construção do saber educacional não pode prescindir da pluralidade de fontes probatórias, uma vez que diferentes tipos de evidência desempenham funções complementares na redução da incerteza inerente ao ato educativo. Nesse sentido, enquanto estudos experimentais de diferentes tipos, revisões sistemáticas e metanálises são cruciais para quantificar efeitos médios e estabelecer generalizações sobre os efeitos de intervenções em larga escala, os relatos e registros sistemáticos de prática documentam os processos, os mecanismos subjacentes e as condições específicas de implementação que determinam a eficácia real de uma intervenção no chão da escola. Cabe, portanto, à expertise profissional refinar a aplicação desses conhecimentos teóricos em sala de aula, transformando dados em práxis.

A necessidade de rigor metodológico torna-se premente na apropriação de Tecnologias Digitais de Informação e Comunicação (TDICs), campo em que a inovação tecnológica costuma preceder a validação pedagógica. Nesse cenário, o uso de videogames atrai interesse crescente tanto pelo apelo lúdico quanto pela sua capacidade arquitetônica de criar cenários simulados. Tais características de simulação e feedback imediato oferecem um suporte de alta fidelidade e baixo risco social, convertendo o erro de um caráter punitivo para uma estrutura construtiva de aprendizagem.

Quando o olhar investigativo se volta ao videogame como objeto de ensino, destaca-se que suas possibilidades de interação e estruturas ramificadas viabilizam a prática deliberada mediante o ensaio repetido de ações. Em títulos com forte eixo narrativo, esse arranjo favorece a resolução de problemas éticos, pois os jogos operam como sistemas fechados e formais que representam subjetivamente um subconjunto da realidade, conforme a definição clássica de \textcite{crawford1984art}. Sob a ótica comportamental, essa natureza de sistema fechado funciona como um ambiente de contingências controladas, no qual o estudante emite respostas e acessa consequências em um espaço de segurança. Desde que acompanhada de mediação sistematizada, a estrutura do jogo garante que as regras atuem na modelagem do repertório, sem produzir os danos típicos de cenários não programados.

Ao mesmo tempo, esses cenários virtuais não operam em isolamento social. A mediação pedagógica atua como variável suplementar à experiência lúdica, ampliando o escopo formativo do jogo e consolidando repertórios estimulados pela interação digital. A interlocução com o campo da comunicação e dos estudos socioculturais dos jogos \cite{mendes2023leitura,nascimento2025retorica} afasta a premissa de que o software seja um ambiente neutro. Sob uma ótica interdisciplinar compatível com as ciências do comportamento, os videogames são compreendidos como artefatos sociotécnicos que codificam práticas culturais por meio de suas estruturas de regras e mecânicas. Desse modo, a intervenção do educador incide sobre a lacuna comunicacional entre a emissão de respostas operantes para o avanço no jogo e a formulação de comportamentos verbais críticos sobre os valores que o sistema exige. Essa articulação é especificamente importante, sobretudo na contemporaneidade, visto que os videogames coexistem com a mobilidade tecnológica, permeando a cultura juvenil de forma onipresente \cite{bogost2011how}.

Não menos importante, o caráter transdisciplinar das TDICs manifesta-se na capacidade de integrar diferentes níveis de análise ambiental e cultural. A compreensão da interação em ambientes digitais exige a articulação de múltiplas áreas de investigação científica. Conforme argumentam  \textcite{todorov2004analise}, a análise de práticas sociais demanda um diálogo contínuo da psicologia com a sociologia e a educação para o estudo efetivo de contingências em nível macro. No caso específico do videogame, essa convergência consolida-se em um ecossistema interconectado \cite{jenkins2006convergence},no qual o game design, a resolução de problemas e a construção de regras mobilizam a arquitetura computacional e os processos comportamentais simultaneamente. As tecnologias configuram-se como cenários de mediação cultural que superam fronteiras disciplinares rígidas.

Diante dessa especificidade midiática, a pesquisa empírica com videogames demanda um delineamento multidimensional. Como destacam \textcite{lieberoth2015mixed}, essa análise contempla diferentes camadas de interação. Essas camadas incluem as experiências \textit{in-game}, referentes aos aspectos intrínsecos do jogo, como narrativa, mecânicas, regras e design de níveis; as experiências \textit{in-room}, focadas nos aspectos situados no espaço físico e nos impactos da presença de outros jogadores nas pessoas presentes; as experiências \textit{in-body}, voltadas ao uso de tecnologias avançadas como biossensores, rastreamento ocular e monitoramento de expressões faciais para fornecer dados fisiológicos adicionais sobre a interação, permitindo uma análise detalhada dos processos comportamentais envolvidos; e as experiências \textit{in-world}, que consideram os elementos da vida pregressa e contextual do indivíduo que refletem em sua experiência de jogabilidade e como o jogo impacta sua vida fora da tela.

Neste recorte de pesquisa, a atenção volta-se primordialmente para a aplicação de critérios \textit{in-game} e \textit{in-room}, uma escolha metodológica justificada pelas impossibilidades técnica, instrumental e temporal para a captura adequada dos demais aspectos em um ambiente escolar padrão, mas que ainda assim permite uma compreensão robusta das dinâmicas de ensino e aprendizagem.

A pertinência dessa delimitação metodológica encontra respaldo empírico nas intervenções direcionadas ao desenvolvimento de competências socioemocionais. A revisão realizada por \textcite{durlak2011impact} indica que práticas educativas que articulam experiência ativa, sequência instrucional lógica e mediação sobre auto-observação têm maior probabilidade de produzir ganhos comportamentais mensuráveis e duradouros. Por isso, estudar os efeitos emergentes em sala de aula mediante o uso de um videogame oferece subsídios pragmáticos indispensáveis para a consolidação da EBE. Tais evidências fundamentam uma transição metodológica rigorosa, substituindo o planejamento pautado em intuição pedagógica por intervenções validadas empiricamente e submetidas à avaliação sistemática.

Autores como \textcite{szczepanik2011concepcao}, ao trabalhar reflexões sobre a filosofia da ciência proposta por \textcite{Bunge1980,hansson2021pseudociencia}, ao definir critérios mais precisos de demarcação entre ciência e pseudociência, fornecem o lastro teórico para essa empreitada. Sob essa ótica, a investigação científica na escola deve indicar mecanismos potenciais para elaborar uma intervenção que, ao definir um objetivo, produza os efeitos esperados para alcançá-lo, forneça condições claras de implementação e destaque as lacunas que demandam instrumentação adicional, operando em um ciclo contínuo de autocorreção e mudança de métodos baseada em dados. Nesse contexto, se por um lado a validade estatística depende da comparação com a média populacional, por outro, o potencial transformador só se concretiza quando tais evidências são adaptadas às idiossincrasias dos sujeitos envolvidos.

Para capturar a complexidade desses fenômenos, o trabalho buscou articular evidências qualitativas provenientes de observação participante, transcrições de rodas de conversa e respostas dissertativas, metodologias amplamente validadas na pesquisa educacional \cite{costa2024likert,kemmis1988action,marques2016observacao,tripp2005pesquisa}.

Entende-se que as evidências qualitativas, longe de serem anedóticas ou secundárias, podem oferecer um entendimento dos processos pragmáticos que ocorrem no cotidiano da sala de aula. Elas fornecem pistas para o aperfeiçoamento de modelos pedagógicos que, quando devidamente manualizados, averiguados e replicados, podem se integrar à Prática Baseada em Evidências (PBE) em Educação, conferindo-lhe a necessária validade para generalização.

Como adendo à discussão sobre a cientificidade na educação, é importante notar que o campo não é isento de controvérsias. Algumas críticas destacam que instituições que promovem a EBE, ao priorizarem práticas sustentadas por evidências quantitativas e por procedimentos metodológicos estritamente padronizados, podem deixar menos espaço para a diversidade das realidades educacionais locais e culturais, promovendo uma ênfase tecnocrática na “eficiência”. Nessa acepção crítica, eficiência é muitas vezes entendida de forma reducionista, como a busca por melhores resultados em testes padronizados com menor investimento humano, financeiro ou temporal, o que pode, em determinadas circunstâncias, desconsiderar efeitos colaterais nocivos, a exemplo do aumento de pressão psicológica sobre docentes e estudantes, da exaustão profissional, além de outros efeitos deletérios no cotidiano institucional, favorecendo a mensuração de resultados em detrimento de abordagens mais amplamente inclusivas e humanizadoras \cite{devechi2022abordagem}.

%confirmar pra mim por favor se a referencia Devechi et al 2022 e essa mesmo q ta na bib

%Está correta sim. Em referências com 3 autores, não precisa usar et al. O autor errou aqui ao usar a expressão.

Trata-se de uma preocupação legítima que merece atenção constante na construção de políticas e práticas embasadas cientificamente, garantindo que a ênfase em evidências não ocorra em detrimento das necessidades contextuais e dos aspectos fundamentais de justiça educacional. De todo modo, no arcabouço teórico que norteou esta pesquisa, a EBE é compreendida através de outros atravessamentos, que priorizam a autonomia intelectual do professor e a integralidade do sujeito aprendiz.

A discussão que se desenrola neste trabalho explora os achados e as reflexões teóricas por uma aproximação à EBE sustentada pelos pilares vinculados a \textcite{slavin1996education,davies1999evidence,young2016conhecimento} que defendem uma educação informada por dados, mas guiada por valores. Tão importante quanto, a pesquisa ancora-se em pressupostos de um fazer científico contextual, falível e autocorrigível, pautado na dúvida sistemática e no ceticismo organizado \cite{hansson2021pseudociencia,szczepanik2011concepcao}. Tal postura reconhece a impossibilidade epistemológica da neutralidade absoluta na ciência, sobretudo nas ciências humanas, apoiando-se, em contrapartida, no ideal de imparcialidade na condução da investigação e na análise dos dados \cite{lacey2011imparcialidade}. Adicionalmente, essa perspectiva científica rigorosa não se deixa sustentar por relativismos epistêmicos ou arbitrariedades metodológicas, defendendo que, embora a redução da incerteza seja uma tarefa árdua, existem critérios objetivos para diferenciar práticas eficazes daquelas inócuas ou prejudiciais \cite{leonardi2023ciencia}.

Aprofundando a relação entre a teoria comportamental e o design de jogos, é fundamental compreender como os videogames operam como tecnologias de ensino sob uma perspectiva analítico-comportamental. \textcite{skinner1974behaviorismo,baum2019compreender} oferecem ferramentas conceituais para entender a aprendizagem por dinâmicas que se desvinculam de processos mentalistas internos e inacessíveis, destrinchando seus pressupostos em mudanças observáveis no comportamento, resultantes da interação do organismo com o ambiente, sem que isso ignore processos internos, história de vida, idiossincrasias e fatores filogenéticos.

Assim, um videogame bem projetado funciona como uma máquina de ensino, discutida por \textcite{skinner1972tecnologia}, mas sofisticada, capaz de fornecer reforço imediato e diferencial para comportamentos específicos, além de superar o isolamento e a rigidez comportamental típicos das máquinas originais.

A retórica procedimental descrita por \textcite{bogost2011how} sustenta a premissa supracitada, tangenciando a semiótica computacional e estabelecendo que a interação em ambientes digitais se consolida prioritariamente pela execução de regras e pela superação da dependência de elementos estritamente visuais ou linguístico-verbais. Nessa dinâmica, o jogador não extrai um sentido abstrato do sistema ao interagir com as mecânicas; ele tem seu repertório comportamental alterado pelas consequências por elas programadas. Essa dimensão semiótica pode ser traduzida, em termos comportamentais, como um arranjo de contingências no qual as regras e a interface do jogo operam como estímulos discriminativos complexos. Tais estímulos sinalizam as condições para modelar o comportamento do jogador, exigindo a emissão de respostas que se mostrem funcionalmente adaptativas para acessar reforçadores ou alterar o ambiente dentro daquele sistema formal.

Essa visão alinha-se também com a noção de “sistemas fechados”, de \textcite{crawford1984art}. Ao isolar variáveis e simplificar a realidade, o jogo permite que o estudante entre em contato com as consequências naturais de suas ações de forma célere e evidente.

Em contingências sociais não programadas, as consequências de uma resposta inadequada, como a ausência de empatia, operam frequentemente sob atraso ou apresentam contornos ambíguos. Em contrapartida, no sistema formal do jogo, as contingências são estruturadas para fornecer consequências claras e imediatas. Essa distinção ambiental cria uma oportunidade específica para o ensino de competências socioemocionais.

A arquitetura do software viabiliza o que a literatura comportamental define como modelagem, permitindo que a aquisição de determinados repertórios sociais ocorra por meio de aproximações sucessivas. O estudante pode testar hipóteses de interação, errar, receber feedback do sistema (seja pelo bloqueio do progresso, pela reação adversa dos personagens ou por alterações na trilha sonora) e ajustar seu comportamento na tentativa seguinte, configurando o ciclo de variação e seleção que fundamenta a aquisição de repertórios operantes.

Ainda sobre a questão da evidência qualitativa e sua validade na EBE, é crucial reiterar o papel da pesquisa-ação e da observação participante como pontes entre a teoria e a prática. \textcite{tripp2005pesquisa,kemmis1988action} argumentam que a pesquisa educacional não deve ser algo feito sobre os professores e os alunos, mas com eles, visando a transformação da realidade. A observação participante \cite{marques2016observacao} permite ao pesquisador acessar o chão da sala de aula, indo além da documentação avaliativa que buscaria descobrir se uma intervenção funcionou, mas como e por que ela funcionou (ou falhou) em um determinado contexto. Isso responde diretamente à demanda de \textcite{davies1999evidence} por evidências que sejam úteis para a tomada de decisão prática.

Saber que videogames aumentam o engajamento em média é útil, mas saber quais tipos de interações verbais e não verbais ocorrem quando um grupo específico de adolescentes joga um título específico fornece um nível de detalhe instrumental que permite a outros educadores replicarem ou adaptarem a experiência com maior chance de sucesso.

\textcite{costa2024likert} reforçam a importância da validade dos instrumentos e métodos qualitativos, argumentando que o rigor na coleta e análise de dados subjetivos (como relatos e percepções) é o que diferencia a pesquisa científica do relato de experiência informal. A triangulação de dados, método central nesta abordagem, serve como um mecanismo de controle de qualidade, permitindo que as limitações de uma fonte de dados (como a desejabilidade social em uma resposta escrita) sejam compensadas pela riqueza de outra (como a observação direta de um comportamento espontâneo).

Este trabalho, portanto, situa-se nesse cruzamento, buscando, através do rigor metodológico e da sensibilidade teórica, compreender as potencialidades reais do videogame como ferramenta para uma formação humana integral.

\section{Metodologia}

O presente estudo constitui um recorte de uma investigação mais ampla, desenvolvida em nível de mestrado, que adotou originalmente uma abordagem de métodos mistos para avaliar o efeito de videogames no ensino formal. Embora a pesquisa matriz tenha também contemplado métricas quantitativas de desempenho e engajamento, este artigo dedica-se exclusivamente à análise qualitativa e funcional das interações verbais e sociais. Tal recorte justifica-se pela necessidade de aprofundar a compreensão das contingências operadas pelo software e pela mediação pedagógica, exigindo uma descrição dos repertórios comportamentais emergentes \cite{skinner1972tecnologia,skinner1974behaviorismo,skinner1978comportamentoverbal}.

Nesse sentido, a investigação adotou uma abordagem qualitativa de natureza exploratória e interpretativa, fundamentada na perspectiva de \textcite{creswell2017research}, para quem esse tipo de pesquisa busca compreender o significado que indivíduos ou grupos atribuem a um problema social. Neste estudo, a complexidade reside na interseção entre tecnologia digital (videogame), desenvolvimento humano (competências socioemocionais) e o ambiente escolar formal. A opção por esse delineamento permitiu surpassar a verificação de frequências de resposta isoladas, focando na interpretação dos processos durante a interação com o objeto de aprendizagem.

A coleta de dados em ambiente natural, a sala de aula de informática e a análise indutiva dos padrões comportamentais possibilitaram capturar as reações socioemocionais e a influência das contingências ambientais imediatas.

O desenho metodológico incorporou elementos da pesquisa-ação, conforme descrito por \textcite{kemmis1988action,tripp2005pesquisa}, uma vez que o pesquisador não atuou como mero observador, mas como um agente mediador que interveio no campo para promover reflexão e ajuste nas práticas pedagógicas durante o curso da investigação. A postura adotada foi a de observação participante \cite{marques2016observacao}, na qual a imersão no contexto e a interação direta com os sujeitos foram utilizadas como ferramentas para acessar os significados atribuídos pelos estudantes às suas experiências. A ênfase nesta vertente qualitativa ampara-se na premissa de que a compreensão dos fenômenos educativos requer um olhar atento às dinâmicas interacionais e aos processos de significação no chão da escola, permitindo uma análise que investiga tanto o resultado quanto o percurso da aprendizagem.

\subsection{Participantes e contexto da pesquisa}

O estudo foi conduzido em uma Escola Técnica Estadual (ETEC) localizada no estado de São Paulo, inserida no contexto do Ensino Médio Integrado ao Técnico (EMIT). A escolha desse cenário justifica-se pela tensão curricular observada na formação técnica, que historicamente prioriza o desenvolvimento de \textit{hard skills} ou competências técnicas em detrimento das \textit{soft skills} ou competências socioemocionais, conforme apontado por \textcite{silvaneto2024corpo,hildebrando2022competencias}. O ambiente físico da intervenção foi o laboratório de informática da instituição, configurado com computadores individuais dispostos em bancadas, o que permitiu tanto a experiência individual de jogo, conhecida como \textit{single-player}, quanto a interação lateral entre os pares.

Os participantes foram 17 estudantes do 2º ano do curso de Desenvolvimento de Sistemas, com idades entre 14 e 16 anos. A seleção da turma obedeceu a critérios pré-definidos pela instituição (permissão de acesso e realização da pesquisa), bem como à pertinência do curso técnico, uma vez que estudantes da área de tecnologia possuem familiaridade prévia com a linguagem dos jogos digitais, o que reduz a curva de aprendizado mecânico e permite focar na análise dos aspectos narrativos e emocionais e, em certa medida, mitigar o distanciamento histórico entre a formação técnica e o desenvolvimento de competências socioemocionais, uma lacuna na formação integral desses estudantes \cite{silvaneto2024corpo}
 
\subsection{Instrumento de intervenção: o videogame como dispositivo pedagógico}


A ferramenta central da intervenção foi o jogo \textit{Lenin – The Lion} \cite{bueno2019lenin}, um RPG (\textit{Role-Playing Game}) eletrônico de autoria do próprio pesquisador. O jogo classifica-se como uma narrativa interativa de aventura e quebra-cabeças, centrada na trajetória de Lenin, um animal antropomorfo que enfrenta exclusão social, bullying e um quadro depressivo severo em sua comunidade. A escolha deste título baseou-se nas características intrínsecas do design, que favorecem a projeção e a reflexão sobre a vida cotidiana, alinhando-se à definição de \textcite{crawford1984art}.

O jogo opera através de uma retórica procedimental \cite{bogost2011how}, na qual as regras e mecânicas de jogo, campo de estudo da ludologia, constroem argumentos persuasivos sobre como o mundo funciona. Em \textit{Lenin – The Lion}, mecânicas como a alteração da paleta de cores de acordo com o estado emocional do protagonista, a alteração da trilha sonora em momentos de crise de pânico e a impossibilidade de realizar certas ações quando o personagem está “deprimido” servem como metáforas operacionais para os sintomas da depressão. A estrutura de escolhas ramificadas permite que os jogadores experimentem consequências éticas e emocionais contingentes às suas decisões, sem os riscos reais do convívio social, criando um simulacro seguro para o ensaio de comportamentos.

A aplicação do jogo visou estimular a emissão de comportamentos correlatos às definições de competências socioemocionais, estipuladas pela Base Nacional Comum Curricular (BNCC) \cite{brasil2018bncc} e discutidas na literatura especializada \cite{durlak2011impact}, tais como empatia, autoconhecimento, regulação emocional e tomada de decisão responsável. Diferentemente de jogos estritamente instrucionais, muitas vezes denominados educational games, que priorizam o conteúdo didático em detrimento da jogabilidade, \textit{Lenin – The Lion} foi utilizado como um objeto cultural capaz de evocar experiências estéticas e éticas complexas, desafiando a concepção de ensino tradicional como norma tácita \cite{jenkins2006convergence}.

\subsection{Procedimentos de coleta de dados e cronograma}

A intervenção desdobrou-se ao longo de quatro encontros quinzenais, com duração aproximada de 1 hora e 30 minutos, correspondendo a duas aulas geminadas. O procedimento foi estruturado para permitir a coleta progressiva de dados, partindo da sensibilização inicial até a auto-observação e análise verbal final.

\subsubsection{Semana 1: sensibilização e diagnóstico}

O primeiro encontro dedicou-se à apresentação dos temas centrais da pesquisa, englobando as competências socioemocionais, o fenômeno do bullying e o potencial dos videogames na educação, além da aplicação de um questionário dissertativo inicial. O objetivo primordial desta etapa foi mapear o conhecimento declarativo prévio e as percepções dos estudantes, estabelecendo uma linha de base para a análise posterior. A observação participante focou nas reações iniciais dos estudantes ao tema e nas expectativas geradas pela proposta de utilizar um videogame em sala de aula, registrando comportamentos de interesse, ceticismo ou curiosidade.

\subsubsection{Semana 2: imersão e primeiras interações}

Os estudantes iniciaram a jogabilidade de \textit{Lenin – The Lion} nos computadores do laboratório. Nesta fase, a coleta de dados concentrou-se na observação direta das interações, utilizando a distinção proposta por \textcite{lieberoth2015mixed} entre interações \textit{in-game}, que ocorrem dentro do universo do jogo, e \textit{in-room}, que acontecem no ambiente físico da sala de aula. Nessa semana, foram registrados, em diário de bordo, comportamentos não verbais, como expressões faciais e gestos, postura e foco atencional, e verbais, incluindo comentários espontâneos, exclamações e pedidos de ajuda técnica ou estratégica, buscando identificar os primeiros indícios de engajamento e apropriação da narrativa em relação aos primeiros dilemas morais apresentados.

\subsubsection{Semana 3: aprofundamento e conflito}

A continuidade do jogo expôs os estudantes a novos dilemas morais, de ordens distintas dos da primeira semana, e a narrativas de exclusão mais intensas dentro da trama. O foco da observação recaiu sobre as estratégias de resolução de problemas adotadas pelos alunos e as justificativas verbalizadas para suas escolhas éticas \textit{in-game}. A mediação do pesquisador intensificou-se neste estágio, utilizando questionamentos socráticos para estimular a verbalização dos processos decisórios e promover a conexão entre os eventos virtuais e as experiências reais dos estudantes. Observou-se como os alunos lidavam com o fracasso, a frustração e as consequências narrativas de suas ações, elementos cruciais para o treino de competências socioemocionais.

\subsubsection{Semana 4: consolidação e auto-observação}

O ciclo encerrou-se com uma Roda de Conversa estruturada e a aplicação do questionário dissertativo final, com as mesmas perguntas do inicial. A Roda de Conversa foi gravada e transcrita para análise posterior, servindo como o principal lócus para a captura de relatos de auto-observação e para a avaliação da apropriação conceitual dos temas trabalhados. O diário de bordo funcionou como instrumento primário de registro contínuo. Nele, o pesquisador estruturou descrições dos eventos, relatos imediatos e interpretações preliminares, fundamentando-se nas recomendações de \textcite{burney2024qualitative} para o uso de métodos qualitativos na pesquisa comportamental.

\subsection{Protocolo de codificação e categorias de análise}

A análise dos dados baseou-se em uma leitura explicitamente analítico-comportamental, buscando operacionalizar conceitos subjetivos em classes de comportamentos observáveis. Para fins de codificação funcional do comportamento verbal, utilizou-se a terminologia clássica da análise do comportamento \cite{skinner1974behaviorismo,skinner1978comportamentoverbal,baum2019compreender}, adaptada para o contexto de interação com mídias digitais e ambientes de aprendizagem, complementarmente definindo metacontingências, a partir de \textcite{todorov2004analise}.

\subsubsection{Mando}

Esta categoria refere-se a enunciados com função solicitadora ou de controle sobre o ambiente ou o ouvinte. No contexto do estudo, foram codificados como mandos os pedidos de ajuda técnica (a exemplo de “como eu abro o menu?”), as solicitações de dicas para resolução de quebra-cabeças (a exemplo de “onde está o diário?”) e as demandas por atenção dos demais participantes e do próprio pesquisador. A frequência e a topografia atrelada à função dos mandos serviram como indicadores do nível de dependência ou autonomia dos estudantes, bem como da natureza colaborativa da tarefa proposta.

\subsubsection{Tato}

O tato descreve a nomeação ou etiquetação de estímulos presentes no ambiente físico ou virtual. Foram classificados nesta categoria os comentários descritivos sobre a narrativa (como “o personagem está triste”), sobre a estética do jogo (como “a música ficou pesada agora”) ou sobre as ações dos personagens (como “eles estão fazendo bullying”). A precisão e a complexidade dos tatos emitidos pelos estudantes funcionaram como indicadores da atenção discriminativa aos elementos sutis do design do jogo e da capacidade de leitura emocional do cenário.

\subsubsection{Intraverbal}

O comportamento intraverbal designa respostas verbais que são evocadas por outros enunciados verbais, sejam do jogo ou de colegas, sem que o objeto falado esteja necessariamente presente fisicamente. Esta categoria englobou as discussões teóricas, as justificativas de escolhas passadas e as conexões feitas entre o jogo e a vida real (como “isso parece com o que aconteceu com meu tio”). A análise dos intraverbais foi crucial para avaliar a capacidade de abstração, a construção de argumentos e a generalização dos conceitos socioemocionais para além do contexto imediato do jogo.

\subsubsection{Relatos de auto-observação}

Além das categorias verbais elementares, o estudo incorporou a categoria de relatos de auto-observação, fundamentada na concepção de autoconhecimento da análise do comportamento \cite{skinner1974behaviorismo,baum2019compreender}. Essa categoria foi utilizada para codificar enunciados em que o falante descreve seu próprio comportamento, monitora suas estratégias ou tateia eventos privados (sentimentos e pensamentos). Exemplos incluem falas nas quais o aluno explicita as variáveis que controlaram seu raciocínio ético (como “você podia ou quebrar o tronco ou não quebrar o tronco. [...] se você não quebrasse, ela tirava o tronco de você. [...] ela, tipo, respeitava você por ter feito aquilo, de algum jeito, entendeu?”) ou seu estado emocional (como “[...] eu fiquei com raiva, eu queria que eles ficassem bravos [também]”).

A presença desses relatos foi interpretada como um indicador de autocontrole e aprofundamento do repertório verbal, marcando a transição da experiência lúdica puramente motora para a aprendizagem consciente e a descrição de contingências.

\subsubsection{Metacontingências: aspectos cerimoniais e tecnológicos}

Para a análise dos fenômenos sociais emergentes na narrativa, adotou-se o referencial teórico proposto por \textcite{todorov2004analise} em sua análise seminal sobre leis e normas sociais. Os autores definem metacontingência como uma relação de contingência entre uma classe de operantes (comportamentos de vários indivíduos) e uma consequência cultural de longo prazo.

Crucialmente, o estudo apropria-se da distinção feita por \textcite{todorov2004analise} entre metacontingências tecnológicas e metacontingências cerimoniais. Enquanto as primeiras envolvem a evolução de práticas visando eficácia ou produção, as cerimoniais são utilizadas por agências de controle (como o estado, ou no caso do jogo, a “sociedade” dos personagens) para garantir a manutenção do status quo, da ordem e da tradição, muitas vezes sendo insensíveis a novas informações ou necessidades individuais. Essa categoria foi a lente utilizada para interpretar a legitimação do preconceito descrita pelos estudantes.

\subsection{Análise Funcional Descritiva das interações}

O protocolo de codificação também contemplou a Análise Funcional Descritiva das interações, fundamentada em \textcite{skinner1974behaviorismo,baum2019compreender}. Para cada episódio comportamental relevante registrado, buscou-se identificar a tríplice contingência que possivelmente o operava.

O antecedente foi identificado como o estímulo que ocasionou o comportamento, podendo ser uma cena específica do jogo, uma pergunta do pesquisador ou um comentário de um colega, sendo o videogame analisado, nesse aspecto, como um arranjo complexo de estímulos discriminativos visuais, auditivos e textuais.

O comportamento foi definido como a resposta observável do estudante, seja ela verbal ou não verbal.

A consequência foi mapeada como o evento que se seguiu ao comportamento e que poderia ter função reforçadora ou punitiva, como o avanço na narrativa, o feedback imediato do jogo, a risada dos colegas ou a validação do comportamento por uma fala.

Essa abordagem permitiu descrever o que os estudantes fizeram e formular hipóteses explicativas sobre o porquê de suas ações, relacionando as mecânicas do jogo \cite{bogost2011how} às respostas emocionais e sociais observadas.

\subsection{Triangulação e validação dos dados}

Para garantir inferências e minimizar vieses interpretativos, adotou-se a estratégia de triangulação de dados \cite{creswell2017research}. A análise integrou informações de três fontes distintas e complementares: a observação direta registrada no diário de bordo, que capturou os comportamentos e contingências em tempo real; o autorrelato oral obtido na Roda de Conversa, cujas transcrições permitiram acesso às interpretações subjetivas e justificativas a posteriori dos estudantes; e o autorrelato escrito proveniente dos questionários dissertativos, que documentaram a evolução do conhecimento declarativo antes e depois da intervenção. A triangulação buscou validar recorrências e mapear padrões consistentes através dessas diferentes fontes.

A título de exemplo, se um comportamento de empatia foi observado durante o jogo (classificado como tato na observação), verificou-se se ele foi verbalizado na discussão subsequente (classificado como intraverbal na Roda de Conversa) e se apareceu na definição conceitual escrita (conhecimento declarativo, no questionário). A convergência entre essas fontes fortaleceu a validade interna do estudo, atendendo aos critérios de rigor propostos para a EBE \cite{davies1999evidence,slavin1996education}. A análise sistemática seguiu um fluxo iterativo que envolveu a leitura flutuante de todo o material, a codificação macro para identificação de temas, a microanálise comportamental e a síntese interpretativa à luz do referencial teórico, assegurando que as conclusões sobre o potencial do videogame fossem sustentadas por evidências empíricas rastreáveis.

\subsection{Aspectos éticos}

A presente investigação pautou-se nos princípios da ética na pesquisa em Ciências Humanas e Sociais, em consonância com as diretrizes da Resolução nº 510/2016 do Conselho Nacional de Saúde (CNS) \cite{brasil2016resolucao510}.

Por tratar-se de uma intervenção pedagógica realizada no contexto natural de sala de aula, a pesquisa obteve a autorização institucional da gestão da unidade escolar para sua realização. A participação dos estudantes foi voluntária e condicionada à assinatura do Termo de Assentimento Livre e Esclarecido (TALE) pelos discentes e do Termo de Consentimento Livre e Esclarecido (TCLE) pelos seus respectivos responsáveis legais. Foi garantido a todos os participantes o direito à desistência a qualquer momento, sem prejuízo, bem como o sigilo absoluto de suas identidades.

\section{Resultados: análise funcional das interações}
A triangulação dos dados coletados via diário de bordo, questionários e, fundamentalmente, a transcrição da Roda de Conversa, permitiu mapear a emergência de operantes verbais complexos. Diferentemente de uma análise de conteúdo tradicional, que buscaria temas, a análise funcional aqui empreendida categorizou as falas dos estudantes conforme a função que exerceram na interação com o ambiente do jogo e com o grupo.

Os dados foram agrupados em três classes comportamentais distintas: (\ref{tatos}) “Tatos complexos e nomeação de eventos privados”; (\ref{intraverbais}) “Intraverbais de tomada de perspectiva”; e (\ref{analise}) “Análise de metacontingências e práticas culturais”.

\subsection{Tatos complexos e nomeação de eventos privados}\label{tatos}

A primeira classe de respostas observada refere-se à capacidade dos estudantes de tatear (descrever) contingências emocionais sutis, indo além da descrição física do cenário. O jogo operou como um estímulo discriminativo (SD) para a verbalização de sentimentos, tanto os projetados nos personagens quanto os sentidos pelos próprios jogadores.

O relato do Estudante A exemplifica a discriminação de uma contingência de depressão familiar, ao descrever a interação do protagonista com a mãe: “Me deu um, tipo, uma tristeza, sabe? Dó, do personagem, sabe? Da própria mãe dele. Aí, eu percebi como que ele se sentia e a própria mãe, sabe? Que um problema pode ocasionar um problema familiar também”. Nesse episódio, o estudante não descreve a cena visualmente, mas tateia a relação funcional entre o comportamento da mãe e o ambiente (“um problema pode ocasionar um problema familiar [...]”).

Outro dado relevante surgiu quando o Estudante B descreveu a contingência de reforçamento positivo ao interagir com uma personagem não hostil (a empregada), contrastando-a com a punição social recebida de outros personagens (os irmãos): “É a cena da empregada, dando um abraço no protagonista, que eu achei muito fofo. [...] Porque eu fiquei com raiva [dos irmãos], eu queria que eles ficassem bravos [também]”. A fala evidencia que o jogo eliciou respostas emocionais (raiva, afeto) que serviram de base para a tomada de decisão do jogador, o qual buscou contracontrolar os personagens aversivos (“queria que eles ficassem bravos”).

\subsection{Intraverbais de tomada de perspectiva}\label{intraverbais}
A categoria de intraverbais englobou respostas verbais controladas pelas falas dos personagens ou pela lógica narrativa, demonstrando a habilidade de tomada de perspectiva (colocar-se no lugar do outro). Os dados sugerem que as mecânicas do jogo estimulavam os estudantes a abandonar uma visão autocentrada para progredir.

Um exemplo sofisticado de intraverbal analítico foi fornecido pelo Estudante C, que identificou uma cadeia de comportamentos na qual a interação direta era impossível, exigindo mediação: “a gente não consegue nem ver o que ele fala, mas dá pra ver que ele tá depressivo [...]. É tudo uma transferência de algo pra alguém que tem um sentimento diferente [...] [da] outra pessoa”. O termo “transferência”, utilizado pelo aluno, denota a compreensão de que os comportamentos no jogo eram mediados por objetos e terceiros, uma leitura funcional da dificuldade de comunicação do personagem depressivo.

Em outro momento, o Estudante B demonstrou um repertório sofisticado de assertividade e aceitação incondicional. Ao dialogar com a personagem que sofria bullying por sua aparência, o aluno optou por não falsear o que se observava (mentir), mas sim por extinguir a função aversiva da característica física. Em vez de negar o atributo, ele o validou como algo natural e isento de culpa: “Nessa parte aí, eu falei que ela tinha sim uma cabeça de martelo, então tudo bem. Que ela realmente tinha uma cabeça grande, mas que não tinha nenhum... problema nenhum. Em ser assim”. Neste episódio verbal, observa-se a formulação de uma regra de conduta inclusiva: a diferença física (o estímulo “cabeça grande”) é descrita, mas o julgamento depreciativo é removido (“não tinha [...] problema nenhum"). O estudante opera uma distinção entre o fato estético e o valor social, promovendo a dignidade da personagem sem recorrer à distorção da materialidade.

\subsection{Análise de metacontingências e práticas culturais}\label{analise}
O nível mais abstrato de análise verbal emergiu quando os estudantes debateram as razões para a exclusão social do protagonista. Neste ponto, os dados corroboram a definição de metacontingência cerimonial de \textcite{todorov2004analise}, entendida como uma prática mantida para preservar a ordem e os costumes de um grupo, mesmo que nociva a indivíduos externos.

O Estudante E ofereceu uma análise precisa dessa dinâmica, identificando que o julgamento moral dos personagens agressores não era individual, mas produto de uma seleção cultural voltada à manutenção do padrão do grupo:

\begin{quote}
    [...] socialmente e na visão delas, elas não estavam erradas, porque, tipo, todo mundo tinha problema com o protagonista, todo mundo. Então mesmo que na nossa visão eles estejam errados, na sociedade que ele estava, elas estavam certas, elas cresceram com ele sendo julgado, sofrendo preconceito etc. Então, mesmo ações que eles acham certas, afetou, né?
\end{quote}

A locução “na sociedade que ele estava, elas estavam certas” demonstra a identificação do ambiente selecionador descrito por \textcite{todorov2004analise}. O aluno percebe que a exclusão do personagem era reforçada positivamente pela sociedade do jogo para garantir a coesão do grupo (o status quo no qual “todos” rejeitam o diferente). O bullying não é identificado como uma maldade intrínseca, mas como prática cerimonial de validação social.

Complementando essa visão contextualista, o Estudante B generalizou o conceito para além do jogo, formulando uma regra ética sobre a multideterminação do comportamento: “Continuando sobre isso, eu acho que ninguém deveria ser realmente culpado 100\% por qualquer coisa. [...] Sempre tem que ter influência sobre tudo que a pessoa viveu, e nem tudo que a pessoa viveu é culpa dela”. Esta verbalização final demonstra a apropriação do conceito na mesma medida que não crava um determinismo ambiental, mas destaca sua influência, afastando-se de explicações mentalistas para uma compreensão das contingências históricas e culturais que moldam as ações humanas, o que, desse modo, constitui emissão de conhecimento declarativo do tema (“saber sobre” o tema).

\section{Discussão: o videogame como tecnologia comportamental na EBE}

A análise dos resultados, à luz da EBE, revela que o videogame \textit{Lenin – The Lion} operou como uma ferramenta de mediação pedagógica eficaz para o desenvolvimento de repertórios analíticos. A ferramenta depende, complementarmente a seus efeitos intrínsecos, das contingências ambientais da sala de aula e da mediação do pesquisador para converter a experiência lúdica em repertório comportamental educativo estruturado.

A discussão sobre as escolhas éticas, como a decisão de preservar ou destruir elementos do cenário, corrobora a teoria de \textcite{crawford1984art} sobre jogos como representações subjetivas da realidade ou, comportamentalmente, como arranjos análogos de contingências, nos quais o software assume a função da comunidade verbal ao programar consequências imediatas para o comportamento de escolha. Nesse ambiente simulado, os estudantes exerceram a "tomada de decisão responsável", descrita na BNCC \cite{brasil2018bncc}, sem os riscos sociais punitivos de uma interação real fracassada.

A retórica procedimental \cite{bogost2011how} materializou-se de forma contundente na análise das metacontingências cerimoniais apresentada na Seção \ref{analise}. As regras do sistema ensinaram aos alunos que o comportamento individual (o bullying contra o personagem) é funcionalmente mantido pela cultura do grupo \cite{todorov2004analise}. A transição epistêmica observada nas falas dos estudantes, que migraram de julgamentos morais simples (a exemplo possível de “eles são maus”) para análises contextuais (como “eles aprenderam assim porque a sociedade reforçou”), constitui um indicador de aprendizagem de competências socioemocionais.

Para a EBE, isso sugere que o videogame pode ser operacionalizado como um laboratório de sociologia aplicada. A evidência de aprendizado reside na capacidade analítica do aluno de identificar as variáveis que controlam o comportamento humano. O relato do Estudante B sobre a aceitação da “cabeça de martelo” ilustra como o jogo permitiu treinar a discriminação entre características físicas e valores morais, promovendo uma ética baseada na observação pragmática e não na esquiva social.

Ademais, os dados sugerem que considerar aspectos vocacionais amplifica o engajamento. Sendo estudantes de um curso técnico em Desenvolvimento de Sistemas, a oportunidade de discutir a qualidade do jogo e aspectos técnicos funcionou como uma operação estabelecedora, aumentando o valor reforçador das interações.

Contudo, o rigor científico exige o reconhecimento das limitações. A presença do pesquisador alterou as contingências da sala de aula, servindo como um estímulo discriminativo (SD) adicional para as falas reflexivas, já que interagia com o público com comportamento intraverbal. De igual modo, o fato de o videogame ter sido desenvolvido pelo próprio pesquisador pode introduzir um viés de autoria na análise dos dados. Essa dupla função de criador do artefato e mediador da intervenção instaura contingências potenciais de desejabilidade social e ameaça da validade externa do estudo.

Os resultados indicam que o videogame serviu como o estímulo eliciador (quando se trata de emoções e sentimentos) e evocador de comportamentos complexos, mas a resposta educativa foi refinada (modelada) pelas interações verbais na Roda de Conversa. Isso reafirma que a tecnologia demanda a intencionalidade pedagógica docente para que o entretenimento dê lugar à formação esperada pelas instituições de ensino, pelas políticas norteadoras expostas na BNCC \cite{brasil2018bncc} e, claro, pelos próprios estudantes.

\section{Conclusão}
A presente investigação, ao articular os pressupostos da EBE com a análise do comportamento e o uso de videogames, permitiu construir um cenário analítico sobre as potencialidades e os desafios da inovação pedagógica no ensino técnico.

Conclui-se que o videogame \textit{Lenin – The Lion} \cite{bueno2019lenin} revelou-se um instrumento mediador de alta eficácia qualitativa. A arquitetura do jogo operou como suporte instrucional, permitindo aos estudantes emitir comportamentos verbais sobre conceitos abstratos que ultrapassavam seu repertório autônomo imediato.

Por meio da Análise Funcional Descritiva, demonstrou-se que a narrativa interativa funcionou como um evocador de respostas classificadas pela comunidade verbal como empatia e auto-observação.

Adicionalmente, a intervenção promoveu o refinamento do conhecimento declarativo dos alunos sobre os temas abordados. Ao identificar metacontingências cerimoniais no ambiente virtual, os participantes transicionaram da emissão de julgamentos morais isolados para uma análise contextual das práticas culturais de exclusão.

A pesquisa levanta novas perguntas fundamentais para a ciência da educação e para a consolidação de uma EBE: como desenhar protocolos de avaliação observacional que capturem a riqueza das interações \textit{in-game} e \textit{in-room} \cite{lieberoth2015mixed} de forma sistematizada e replicável em larga escala? De que maneira a sistematização dessas práticas mediadas pode mitigar o viés da novidade e da presença do pesquisador, permitindo que o videogame se integre à rotina escolar de forma sustentável e autônoma pelos professores \cite{slavin1996education}, desviando-se de um protocolo enrijecido, na mesma medida que orienta a tomada de decisão?

Em suma, o estudo aponta que a operacionalização das competências socioemocionais via videogames é viável e desejável, desde que sustentada por uma intencionalidade pedagógica clara e sujeita a um escrutínio contínuo. A EBE, neste contexto, serve para oferecer um método de autocorreção que permite ao educador distinguir entre o que é meramente lúdico, intrínseco ao videogame, e o que é efetivamente formativo. O videogame na escola, portanto, deve ser encarado menos como um fim em si mesmo e mais como um complexo arranjo de contingências capaz de humanizar, contextualizar e amplificar a experiência de aprendizagem técnica.


\printbibliography\label{sec-bib}
% if the text is not in Portuguese, it might be necessary to use the code below instead to print the correct ABNT abbreviations [s.n.], [s.l.]
%\begin{portuguese}
%\printbibliography[title={Bibliography}]
%\end{portuguese}


%full list: conceptualization,datacuration,formalanalysis,funding,investigation,methodology,projadm,resources,software,supervision,validation,visualization,writing,review
\begin{contributors}[sec-contributors]
\authorcontribution{João Antonio Lemes Bueno}[conceptualization,datacuration,formalanalysis,investigation,methodology,software,visualization,writing,review]
\authorcontribution{Adriana Aparecida de Lima Terçariol}[projadm,supervision,review]
\end{contributors}

\begin{dataavailability}
\txtdataavailability{datanotav} % options: dataavailable, dataonly, databody, datanotav, nodata
\end{dataavailability}

\begin{funding}
O presente trabalho foi realizado com apoio da Coordenação de Aperfeiçoamento de Pessoal de Nível Superior - Brasil (CAPES) - Código de Financiamento 001 (Programa PROSUP/Taxa).
\end{funding}


\end{document}

